\section*{Capítulo \ref{ch:b3:developmental-genetics} --- Genética del desarrollo y evolución de la forma}

\begin{solution}{exo:b3:developmental-genetics:1}
Efecto materno: \emph{bicoid}; los embriones de madres mutantes carecen
de cabeza y de tórax. Gap: \emph{Krüppel}; falta un bloque contiguo de
segmentos. Regla par: \emph{fushi tarazu} (o \emph{even-skipped});
falta un segmento de cada dos. Polaridad segmentaria: \emph{hedgehog}
(o \emph{engrailed}); parte de cada segmento queda sustituida por una
imagen especular del resto.
\end{solution}

\begin{solution}{exo:b3:developmental-genetics:2}
Un gen cuyo producto deposita la madre en el huevo, de modo que el
fenotipo del embrión sigue al genotipo de ella. Sus embriones carecen
de cabeza y de tórax y llevan estructuras posteriores en ambos
extremos; mueren. El alelo silvestre paterno está en el cigoto, pero el
ARNm de \emph{bicoid} lo fabrican durante la ovogénesis las células
nodriza de la madre y queda anclado en el polo antes de la fecundación;
la copia propia del cigoto se transcribiría horas más tarde, cuando el
tiempo del gradiente ya ha pasado.
\end{solution}

\begin{solution}{exo:b3:developmental-genetics:3}
\omterm{def:b3:developmental-genetics:hox}{Colinealidad}: el orden de los \omterm{def:b3:developmental-genetics:hox}{genes Hox} a lo largo del cromosoma
coincide con el orden de sus dominios de expresión a lo largo del
cuerpo. Prevalencia posterior: allí donde se expresan varios, el más
posterior fija la identidad. Un ratón sin \emph{Hoxc8} muestra una
transformación anterior ---su primera vértebra lumbar toma la identidad
de una torácica y lleva una costilla. \emph{Antennapedia} en la cabeza
convierte las antenas en patas, la identidad del segundo segmento
torácico.
\end{solution}

\begin{solution}{exo:b3:developmental-genetics:4}
Un labio dorsal injertado indujo un segundo eje completo hecho sobre
todo de células del huésped: el injerto instruyó a sus vecinas ---un
\omterm{def:b3:developmental-genetics:organiser}{organizador}. Sus moléculas son antagonistas secretados de la BMP
(cordina, noguina, folistatina). El ectodermo se vuelve neural cuando
falta la señalización BMP y epidérmico cuando la BMP está presente, de
modo que lo neural es lo que hace el ectodermo cuando no se le dice
nada: el \omterm{def:b3:developmental-genetics:organiser}{organizador} actúa silenciando una instrucción.
\end{solution}

\begin{solution}{exo:b3:developmental-genetics:5}
$k = \ln 2/(35\times 60) = \qty{3.3e-4}{s^{-1}}$; $\lambda = \sqrt{4/
3.3\times 10^{-4}} = \qty{110}{\micro m}$; razón $\mathrm{e}^{250/110}
= \mathrm{e}^{2.27} = 9.7$.
\end{solution}

\begin{solution}{exo:b3:developmental-genetics:6}
Cuatro copias duplican $c_{0}$: el límite retrocede $\lambda\ln 2 =
\qty{69}{\micro m}$ hasta \qty{309}{\micro m}; una sola copia lo reduce
a la mitad: avanza \qty{69}{\micro m} hasta \qty{171}{\micro m}. Cada
desplazamiento es el \qty{14}{\%} de la longitud del huevo ---mayor que
lo observado, lo que ya es indicio de mecanismos que amortiguan la
dependencia de la dosis.
\end{solution}

\begin{solution}{exo:b3:developmental-genetics:7}
$\delta x = \lambda\,\delta c/c$: para \qty{8}{\micro m}, $\delta c/c =
8/100 = \qty{8}{\%}$. Un recuento de Poisson $N$ tiene un error
relativo $1/\sqrt{N}$, de modo que hay que contar
$N = 1/0.08^{2} \approx 160$ moléculas ---un pequeño porcentaje de las
moléculas presentes en un núcleo, o un único recuento integrado en el
tiempo.
\end{solution}

\begin{solution}{exo:b3:developmental-genetics:8}
$k = \ln 2/35 = \qty{0.0198}{min^{-1}}$: $1 - \mathrm{e}^{-kt}$ da
$0.70$ a los \qty{60}{min}, $0.83$ a los \qty{90}{min}, $0.91$ a los
\qty{120}{min} y $0.95$ a los \qty{150}{min}. El gradiente sigue
subiendo un pequeño porcentaje mientras se lee, lo que desplazaría un
límite $\lambda\ln(1/0.9) \approx \qty{10}{\micro m}$, alrededor de un
núcleo, a lo largo de la ventana de lectura: el embrión ha de leer a un
tiempo fijo o emplear una lectura insensible al nivel global.
\end{solution}

\begin{solution}{exo:b3:developmental-genetics:9}
Una segunda ZAP da una duplicación en espejo, 4-3-2-2-3-4: las células
anteriores ven ahora un Shh alto y se vuelven dedos posteriores. Una
bolita que libera poco Shh da una duplicación parcial ---un dedo 2
anterior de más, pongamos 2-2-3-4--- porque las células anteriores solo
ven la concentración baja que especifica el dedo 2.
\end{solution}

\begin{solution}{exo:b3:developmental-genetics:10}
Solución particular $s/k$; soluciones homogéneas $\cosh(x/\lambda)$ y
$\sinh(x/\lambda)$; la ausencia de flujo en $0$ elimina el $\sinh$;
$c(L) = 0$ fija la amplitud:
$c(x) = (s/k)\bigl[1 - \cosh(x/\lambda)/\cosh(L/\lambda)\bigr]$. El
perfil es una meseta cerca de $x = 0$ que cae a cero en una capa límite
de anchura $\lambda$ junto al sumidero ---no una exponencial desde un
punto, sino un campo llano con un borde; la concentración más alta es
la más lejana al sumidero. Para $L \ll \lambda$, $\cosh u \approx 1 +
u^{2}/2$ y $c \approx s(L^{2} - x^{2})/(2D)$: una parábola independiente
de $k$, ya que las moléculas alcanzan el sumidero antes de poder
degradarse.
\end{solution}

\begin{solution}{exo:b3:developmental-genetics:11}
$x_{1} = \lambda\ln(1/0.3) = 1.20\lambda$ y $x_{2} = \lambda\ln(1/0.08)
= 2.53\lambda$: la primera región mide $1.20\lambda$ de ancho y la
segunda $1.32\lambda$, ambas independientes de $c_{0}$, que las
desplaza juntas. La tercera región, $L - 2.53\lambda$, absorbe toda la
variación de la longitud del huevo: en un huevo de \qty{450}{\micro m}
con $\lambda = \qty{100}{\micro m}$ el abdomen mide
\qty{197}{\micro m}, y en uno de \qty{550}{\micro m},
\qty{297}{\micro m} ---el patrón no está escalado. Mecanismos: un
segundo gradiente desde el polo posterior (Nanos, Caudal o el sistema
terminal) leído junto con Bicoid, de modo que las posiciones las fije
la razón; o una molécula ``expansora'' fabricada allí donde el
morfógeno es bajo que alargue $\lambda$ hasta que el gradiente llene el
huevo.
\end{solution}

\begin{solution}{exo:b3:developmental-genetics:12}
Un cambio codificante altera la proteína allí donde actúe:
\emph{Pitx1} construye además la \omterm{def:b3:endocrinology:axes}{hipófisis} y la mandíbula, y \emph{\omterm{def:b3:developmental-genetics:organiser}{sonic
hedgehog}} da además su patrón al tubo neural, al intestino y a los
dientes, de modo que una proteína rota es letal o pleiotrópicamente
incapacitante. Un potenciador dirige la expresión en un tejido y en un
momento; suprimirlo retira la estructura que ese tejido habría
construido y deja intacto todo otro uso de la proteína. Los
potenciadores son modulares y a menudo en parte redundantes, de modo
que los pasos intermedios son viables, y la misma ruta ---un
interruptor perdido--- se ha tomado de forma independiente en muchos
lagos de espinosos.
\end{solution}

\begin{solution}{pb:b3:developmental-genetics:1}
\textbf{1.} $k = \ln 2/(35\times 60) = \qty{3.3e-4}{s^{-1}}$; $1/k =
\qty{3030}{s} = \qty{50}{min}$.
\textbf{2.} $\lambda = \sqrt{4/3.3\times 10^{-4}} = \qty{110}{\micro m}$;
el huevo mide $4.5\lambda$.
\textbf{3.} $\mathrm{e}^{4.5} \approx 90$ de un polo al otro;
$\mathrm{e}^{8.5/110} = 1.08$, un cambio del \qty{8}{\%} por núcleo.
\textbf{4.} $J = 50\times\sqrt{4\times 3.3\times 10^{-4}} = 50\times
0.036 = \qty{1.8}{nM\,\micro m\,s^{-1}}$. Un nM son $0.6$ moléculas
por \unit{\micro m^3}, de modo que son unas $1.1$ moléculas por
\unit{\micro m^2} y segundo.
\textbf{5.} Volumen nuclear $\tfrac{4}{3}\pi\times 27 = \qty{113}{\micro
m^3}$. A \qty{50}{nM}, $30$ moléculas por \unit{\micro m^3}:
\num{3400} moléculas. En el centro del huevo, $c = 50\,\mathrm{e}^{-2.27} =
\qty{5.2}{nM}$: unas \num{350} moléculas.
\textbf{6.} $kt = 1.19$ a los \qty{60}{min}: $0.70$; $kt = 2.38$ a los
\qty{120}{min}: $0.91$. El gradiente está a menos del \qty{10}{\%} del
estacionario cuando se lee ---y sigue subiendo lo bastante como para
mover un límite $\lambda\ln(1/0.91) \approx \qty{10}{\micro m}$, un
núcleo.
\textbf{7.} $x_{1} = 110\ln(1/0.3) = \qty{132}{\micro m}$
(\qty{26}{\%}); $x_{2} = 110\ln(1/0.08) = \qty{278}{\micro m}$
(\qty{56}{\%}).
\textbf{8.} Cuatro copias: $c_{0} = \qty{100}{nM}$; ambos límites
retroceden $\lambda\ln 2 = \qty{76}{\micro m}$, hasta $208$ y
\qty{354}{\micro m}. Seis copias: $\lambda\ln 3 = \qty{121}{\micro m}$,
hasta $253$ y \qty{399}{\micro m}.
\textbf{9.} Una copia: avanzan \qty{76}{\micro m}, hasta $56$ y
\qty{202}{\micro m}. La región de la cabeza encogió de $132$ a
\qty{56}{\micro m}; el tórax mantuvo su anchura (\qty{146}{\micro m}) y
se adelantó; el abdomen creció de $222$ a \qty{298}{\micro m}.
\textbf{10.} $\delta x = 110\times 0.1 = \qty{11}{\micro m}$, $1.3$
espaciados nucleares.
\textbf{11.} Un núcleo: $\delta c/c = 8.5/110 = \qty{7.7}{\%}$. Desde el
\qty{10}{\%}, $n = (10/7.7)^{2} \approx 1.7$: bastan dos lecturas
independientes (o el promedio de dos núcleos vecinos).
\textbf{12.} $1/\sqrt{350} = \qty{5.3}{\%}$: un solo recuento
instantáneo es ya más preciso que el \qty{10}{\%} supuesto, de modo que
el ruido observado no es solo ruido de recuento, sino que procede de la
cinética de unión y de la maquinaria de lectura ---y promediar en el
tiempo lo mejora.
\textbf{13.} $278/450 = \qty{62}{\%}$ y $278/550 = \qty{51}{\%}$ de la
longitud del huevo: una discrepancia del \qty{11}{\%}, seis núcleos. El
patrón no escala con un $\lambda$ fijo.
\textbf{14.} Si $D \propto L^{2}$ con $k$ fijo, $\lambda \propto L$ y
$x_{T}/L = (\lambda/L)\ln(c_{0}/T)$ es el mismo en todo huevo: escalado
perfecto. Pero $D$ es una propiedad de la molécula y del citoplasma, no
del tamaño del huevo, de modo que así planteado es inverosímil; en los
embriones reales el escalado es parcial y procede de otras lecturas.
\textbf{15.} $c^{5}/(c^{5} + T^{5}) = 0.1$ en $c/T = 9^{-1/5} = 0.64$ y
$0.9$ en $c/T = 9^{1/5} = 1.55$: un factor $2.4$ en concentración,
$\Delta x = 110\ln 2.4 = \qty{97}{\micro m}$ ---once núcleos, mucho más
ancho que el límite observado de dos o tres; la lectura del gradiente
por sí sola no es lo bastante afilada.
\textbf{16.} La inclinación asigna una concentración dada a encendido o
apagado con más decisión, de modo que la región de transición es más
estrecha; pero un error del \qty{10}{\%} en la concentración sigue
moviendo el punto donde se cruza el umbral en $\lambda\,\delta c/c$: la
respuesta puede ser un escalón perfecto y el escalón caer aun así en el
lugar equivocado.
\textbf{17.} Si ambos estuvieran encendidos, cada uno estaría
reprimiendo al otro, de modo que al menos uno acabaría apagado; con
represión cooperativa los únicos estados estables son uno encendido y
otro apagado, y un límite entre ellos es un interruptor, no una
pendiente.
\textbf{18.} $\mathrm{d}\ln\lambda = \tfrac{1}{2}(0.02 - 0.06) = -0.02$
por grado: $\lambda$ cae un \qty{2}{\%} por grado, un \qty{10}{\%} para
\qty{5}{\celsius}, hasta \qty{99}{\micro m}; el límite de
\emph{hunchback} pasa de $278$ a \qty{250}{\micro m}, tres núcleos más
adelante. Las moscas se desarrollan entre \qty{18}{} y
\qty{29}{\celsius} con proporciones normales, de modo que algo
compensa ---una fuente o una lectura que se desplace con la temperatura
en sentido contrario.
\textbf{19.} $2^{8} = 256$ códigos; se usan catorce. Con prevalencia
posterior solo importa el gen encendido más posterior, de modo que hay
a lo sumo nueve identidades distintas (ocho genes, o ninguno): el
código es una jerarquía, no una palabra binaria.
\textbf{20.} El tercer segmento torácico se convierte en un segundo:
\emph{Ubx} reprimía el programa del ala que un segmento torácico
ejecuta por omisión, y retirar el represor libera el estado ancestral
---cuatro alas, como en los antepasados de cuatro alas de la mosca.
\textbf{21.} La región cervical del ganso llega a la vértebra 22: un
cuello largo de muchas vértebras, frente a las siete del ratón. El gen
es el mismo; lo que se desplazó hacia atrás fue el límite anterior de
su expresión ---un cambio regulador en dónde se enciende un gen Hox.
\textbf{22.} Cada borde es ahora una fuente; el Shh es máximo en ambos
extremos (dedo 4), cae hacia el centro (3, luego 2) y nunca alcanza el
nivel bajo que especifica el dedo 1, de modo que en el medio hay dos
dedos 2 espalda contra espalda: 4-3-2-2-3-4.
\textbf{23.} \emph{Pax6} es un regulador maestro conservado que puede
disparar el programa ocular que lleve el huésped: la mosca construyó un
ojo de mosca, de modo que el gen de ratón no lleva información de ojo
de ratón, sino solo la orden ``constrúyase aquí un ojo''. Muestra una
\omterm{def:b3:developmental-genetics:evodevo}{homología profunda} de la red génica, no una homología de los órganos:
el ojo compuesto y el ojo en cámara evolucionaron por separado a partir
de un antepasado que tenía un parche de fotorreceptores con
\emph{Pax6}.
\textbf{24.} Cinco dedos a lo ancho de \qty{2}{mm} son una longitud de
onda de \qty{0.4}{mm}; seis exigen \qty{0.33}{mm}, un acortamiento de
una sexta parte. Como $\lambda \propto \sqrt{D_{\text{inh}}/k_{\text{inh}}}$,
elévese un \qty{44}{\%} la degradación del inhibidor o redúzcase un
\qty{31}{\%} su difusión; en la placa de la mano lo consigue bajar la
dosis de los \omterm{def:b3:developmental-genetics:hox}{genes Hox} posteriores, y los ratones mutantes tienen seis,
siete o más dedos finos.
\textbf{25.} $\lambda \approx \qty{110}{\micro m}$; una duplicación de
la fuente desplaza todos los límites \qty{76}{\micro m}; un ruido del
\qty{10}{\%} en la concentración son \qty{11}{\micro m}, alrededor de un
núcleo, de error posicional.
\end{solution}
